% Options for packages loaded elsewhere
\PassOptionsToPackage{unicode}{hyperref}
\PassOptionsToPackage{hyphens}{url}
\PassOptionsToPackage{dvipsnames,svgnames,x11names}{xcolor}
\documentclass[
  10pt,
  fontset=fandol]{ctexart}
\usepackage{xcolor}
\usepackage[margin=16mm]{geometry}
\usepackage{amsmath,amssymb}
\setcounter{secnumdepth}{-\maxdimen} % remove section numbering
\usepackage{iftex}
\ifPDFTeX
  \usepackage[T1]{fontenc}
  \usepackage[utf8]{inputenc}
  \usepackage{textcomp} % provide euro and other symbols
\else % if luatex or xetex
  \usepackage{unicode-math} % this also loads fontspec
  \defaultfontfeatures{Scale=MatchLowercase}
  \defaultfontfeatures[\rmfamily]{Ligatures=TeX,Scale=1}
\fi
\usepackage{lmodern}
\ifPDFTeX\else
  % xetex/luatex font selection
\fi
% Use upquote if available, for straight quotes in verbatim environments
\IfFileExists{upquote.sty}{\usepackage{upquote}}{}
\IfFileExists{microtype.sty}{% use microtype if available
  \usepackage[]{microtype}
  \UseMicrotypeSet[protrusion]{basicmath} % disable protrusion for tt fonts
}{}
\makeatletter
\@ifundefined{KOMAClassName}{% if non-KOMA class
  \IfFileExists{parskip.sty}{%
    \usepackage{parskip}
  }{% else
    \setlength{\parindent}{0pt}
    \setlength{\parskip}{6pt plus 2pt minus 1pt}}
}{% if KOMA class
  \KOMAoptions{parskip=half}}
\makeatother
\usepackage{graphicx}
\makeatletter
\newsavebox\pandoc@box
\newcommand*\pandocbounded[1]{% scales image to fit in text height/width
  \sbox\pandoc@box{#1}%
  \Gscale@div\@tempa{\textheight}{\dimexpr\ht\pandoc@box+\dp\pandoc@box\relax}%
  \Gscale@div\@tempb{\linewidth}{\wd\pandoc@box}%
  \ifdim\@tempb\p@<\@tempa\p@\let\@tempa\@tempb\fi% select the smaller of both
  \ifdim\@tempa\p@<\p@\scalebox{\@tempa}{\usebox\pandoc@box}%
  \else\usebox{\pandoc@box}%
  \fi%
}
% Set default figure placement to htbp
\def\fps@figure{htbp}
\makeatother
\setlength{\emergencystretch}{3em} % prevent overfull lines
\providecommand{\tightlist}{%
  \setlength{\itemsep}{0pt}\setlength{\parskip}{0pt}}
\usepackage{amsmath,amssymb,mathtools,needspace}
\xeCJKDeclareCharClass{CJK}{"2032,"0400->"04FF,"2160->"216B,"2460->"2473,"25A0,"25C7,"25CB,"25CF}
\IfFontExistsTF{Arial Unicode MS}{\xeCJKsetup{AutoFallBack=true}\setCJKfallbackfamilyfont{\CJKrmdefault}{Arial Unicode MS}}{}
\setcounter{secnumdepth}{0}
\setlength{\emergencystretch}{2em}
\allowdisplaybreaks
\usepackage{bookmark}
\IfFileExists{xurl.sty}{\usepackage{xurl}}{} % add URL line breaks if available
\urlstyle{same}
\hypersetup{
  pdftitle={Lagrange插值},
  colorlinks=true,
  linkcolor={Maroon},
  filecolor={Maroon},
  citecolor={Blue},
  urlcolor={Blue},
  pdfcreator={LaTeX via pandoc}}

\title{Lagrange插值}
\author{}
\date{}

\begin{document}
\maketitle

\subsection{2.2 Lagrange 插值}\label{lagrange-ux63d2ux503c}

\subsubsection{2.2.1 插值多项式的存在唯一性}\label{ux63d2ux503cux591aux9879ux5f0fux7684ux5b58ux5728ux552fux4e00ux6027}

设 \(P(x)\) 是形如式（2.1.2）的插值多项式，用 \(H_n\) 代表所有次数不超过 \(n\) 的多项式集合，于是 \(P(x)\in H_n\)（符号 \(\in\) 表示属于）。所谓插值多项式 \(P(x)\) 存在唯一，就是指在集合 \(H_n\) 中有且只有一个 \(P(x)\) 满足式（2.1.1）。由式（2.1.1）可得

\[
\begin{cases}
a_0+a_1x_0+\cdots+a_nx_0^n=y_0,\\
a_0+a_1x_1+\cdots+a_nx_1^n=y_1,\\
\qquad\qquad\vdots\\
a_0+a_1x_n+\cdots+a_nx_n^n=y_n.
\end{cases}
\tag{2.2.1}
\]

这是一个关于 \(a_0,a_1,\cdots,a_n\) 的 \(n+1\) 元线性方程组。要证明插值多项式的存在唯一，只要证明方程组（2.2.1）的解存在唯一，也就是证明方程组（2.2.1）的系数行列式

\[
V_n(x_0,x_1,\cdots,x_n)=
\begin{vmatrix}
1 & x_0 & x_0^2 & \cdots & x_0^n\\
1 & x_1 & x_1^2 & \cdots & x_1^n\\
\vdots & \vdots & \vdots & & \vdots\\
1 & x_n & x_n^2 & \cdots & x_n^n
\end{vmatrix}
\tag{2.2.2}
\]

不为零，其中 \(V_n(x_0,x_1,\cdots,x_n)\) 称为 Vandermonde 行列式。利用行列式性质可得

\[
V_n(x_0,x_1,\cdots,x_n)
=\prod_{i=1}^{n}\prod_{j=0}^{i-1}(x_i-x_j),
\]

由于 \(i\ne j\) 时 \(x_i\ne x_j\)，上式所有因子 \(x_i-x_j\ne0\)，于是 \(V_n(x_0,x_1,\cdots,x_n)\ne0\)，故方程组（2.2.1）存在唯一的一组解 \(a_0,a_1,\cdots,a_n\)。以上论述可写成如下定理。

\textbf{定理 2.1} 满足条件式（2.1.1）的插值多项式（2.1.2）存在唯一。

\href{../../source-images/pdf-029.jpeg}{原页：书页 16／PDF 29}

\subsubsection{2.2.2 线性插值与抛物插值}\label{ux7ebfux6027ux63d2ux503cux4e0eux629bux7269ux63d2ux503c}

从定理 2.1 的证明可看到，要求插值多项式 \(P(x)\)，可以通过求方程组（2.2.1）的解 \(a_0,a_1,\cdots,a_n\) 得到。但这样做不但计算复杂，且难以得到 \(P(x)\) 的简单表达式。为了求得便于使用的简单的插值多项式 \(P(x)\)，先讨论 \(n=1\) 的情形。

假定已知区间 \([x_k,x_{k+1}]\) 的端点处的函数值 \(y_k=f(x_k)\)，\(y_{k+1}=f(x_{k+1})\)，要求线性插值多项式 \(L_1(x)\)，使它满足

\[
L_1(x_k)=y_k,\qquad L_1(x_{k+1})=y_{k+1}.
\]

\(y=L_1(x)\) 的几何意义就是通过两点 \((x_k,y_k)\) 和 \((x_{k+1},y_{k+1})\) 的直线，如图 2.2 所示。\(L_1(x)\) 的表达式可由几何意义直接给出，即

\[
L_1(x)=y_k+\frac{y_{k+1}-y_k}{x_{k+1}-x_k}(x-x_k)
\quad\text{（点斜式）},
\tag{2.2.3}
\]

\[
L_1(x)=\frac{x_{k+1}-x}{x_{k+1}-x_k}y_k
+\frac{x-x_k}{x_{k+1}-x_k}y_{k+1}
\quad\text{（两点式）}.
\tag{2.2.4}
\]

\includegraphics[width=0.85\linewidth,height=\textheight,keepaspectratio,alt={图 2.2（原图裁切）}]{../assets/ch02a/figure-2.2.png}

图 2.2（原图）：直线 \(y=L_1(x)\) 与曲线 \(y=f(x)\) 在两节点相交。全部文字标签：\(y\)、\(x\)、\(O\)、\(x_k\)、\(x_{k+1}\)、\(y_k\)、\(y_{k+1}\)、\(y=L_1(x)\)、\(y=f(x)\)。

由两点式方程看出，\(L_1(x)\) 是由两个线性函数

\[
l_k(x)=\frac{x-x_{k+1}}{x_k-x_{k+1}},
\qquad
l_{k+1}(x)=\frac{x-x_k}{x_{k+1}-x_k}
\]

的线性组合得到的，其系数分别为 \(y_k\) 及 \(y_{k+1}\)，即

\[
L_1(x)=y_kl_k(x)+y_{k+1}l_{k+1}(x).
\tag{2.2.5}
\]

显然，\(l_k(x)\) 及 \(l_{k+1}(x)\) 也是线性插值多项式，在节点 \(x_k\) 及 \(x_{k+1}\) 上满足条件

\[
l_k(x_k)=1,\quad l_k(x_{k+1})=0,\quad
l_{k+1}(x_k)=0,\quad l_{k+1}(x_{k+1})=1.
\]

称函数 \(l_k(x)\) 及 \(l_{k+1}(x)\) 为\textbf{一次插值基函数}或\textbf{线性插值基函数}，它们的图形分别如图 2.3(a)、(b) 所示。

\includegraphics[width=0.85\linewidth,height=\textheight,keepaspectratio,alt={图 2.3（原图裁切）}]{../assets/ch02a/figure-2.3.png}

图 2.3（原图）：(a) \(l_k(x)\) 在 \(x_k\) 处为 \(1\)，在 \(x_{k+1}\) 处为 \(0\)；(b) \(l_{k+1}(x)\) 在 \(x_k\) 处为 \(0\)，在 \(x_{k+1}\) 处为 \(1\)。两图全部文字标签：\(y\)、\(x\)、\(O\)、\(1\)、\(x_k\)、\(x_{k+1}\)、\(l_k(x)\)、\(l_{k+1}(x)\)、(a)、(b)。

下面讨论 \(n=2\) 的情况。假定插值节点为 \(x_{k-1},x_k,x_{k+1}\)，要求二次插值多项式 \(L_2(x)\)，使它满足

\begin{quote}
页末插值条件续 PDF 30。
\end{quote}

\href{../../source-images/pdf-030.jpeg}{原页：书页 17／PDF 30}

\begin{quote}
页首承接 PDF 29 的二次插值条件；末式展开续 PDF 31。
\end{quote}

\[
L_2(x_j)=y_j\qquad(j=k-1,k,k+1).
\]

我们知道，\(y=L_2(x)\) 在几何上就是通过三点 \((x_{k-1},y_{k-1})\)、\((x_k,y_k)\)、\((x_{k+1},y_{k+1})\) 的抛物线。为了求出 \(L_2(x)\) 的表达式，可采用基函数方法，此时基函数 \(l_{k-1}(x),l_k(x)\) 及 \(l_{k+1}(x)\) 是二次函数，且在节点上满足条件

\[
\begin{cases}
l_{k-1}(x_{k-1})=1,\quad l_{k-1}(x_j)=0\quad(j=k,k+1),\\
l_k(x_k)=1,\quad l_k(x_j)=0\quad(j=k-1,k+1),\\
l_{k+1}(x_{k+1})=1,\quad l_{k+1}(x_j)=0\quad(j=k-1,k).
\end{cases}
\tag{2.2.6}
\]

满足条件式（2.2.6）的插值基函数是很容易求出的，例如求 \(l_{k-1}(x)\)，因它有两个零点 \(x_k\) 及 \(x_{k+1}\)，故可表示为

\[
l_{k-1}(x)=A(x-x_k)(x-x_{k+1}),
\]

其中 \(A\) 为待定系数，可由条件 \(l_{k-1}(x_{k-1})=1\) 定出，即

\[
A=\frac{1}{(x_{k-1}-x_k)(x_{k-1}-x_{k+1})},
\]

故

\[
l_{k-1}(x)=
\frac{(x-x_k)(x-x_{k+1})}
{(x_{k-1}-x_k)(x_{k-1}-x_{k+1})}.
\]

同理，

\[
l_k(x)=
\frac{(x-x_{k-1})(x-x_{k+1})}
{(x_k-x_{k-1})(x_k-x_{k+1})},
\qquad
l_{k+1}(x)=
\frac{(x-x_{k-1})(x-x_k)}
{(x_{k+1}-x_{k-1})(x_{k+1}-x_k)}.
\]

函数 \(l_{k-1}(x),l_k(x),l_{k+1}(x)\) 称为\textbf{二次插值基函数}或\textbf{抛物插值基函数}，它们在区间 \([x_{k-1},x_{k+1}]\) 上的图形分别如图 2.4(a)、(b)、(c) 所示。

\includegraphics[width=0.85\linewidth,height=\textheight,keepaspectratio,alt={图 2.4（原图裁切）}]{../assets/ch02a/figure-2.4.png}

图 2.4（原图）：(a) \(l_{k-1}(x)\)，(b) \(l_k(x)\)，(c) \(l_{k+1}(x)\)，分别在对应节点取 \(1\)、其他两个节点取 \(0\)；左右基函数在部分区间为负。三图全部文字标签：\(y\)、\(x\)、\(O\)、\(1\)、\(x_{k-1}\)、\(x_k\)、\(x_{k+1}\)、\(l_{k-1}(x)\)、\(l_k(x)\)、\(l_{k+1}(x)\)、(a)、(b)、(c)。

利用二次插值基函数 \(l_{k-1}(x),l_k(x),l_{k+1}(x)\)，立即可得到二次插值多项式

\[
L_2(x)=y_{k-1}l_{k-1}(x)+y_kl_k(x)+y_{k+1}l_{k+1}(x),
\tag{2.2.7}
\]

显然，它满足条件

\[
L_2(x_j)=y_j\qquad(j=k-1,k,k+1).
\]

将上面求得的 \(l_{k-1}(x),l_k(x),l_{k+1}(x)\) 代入式（2.2.7），得

\href{../../source-images/pdf-031.jpeg}{原页：书页 18／PDF 31}

\begin{quote}
页首续 PDF 30 的式（2.2.7）展开；页末公式续 PDF 32。
\end{quote}

\[
\begin{aligned}
L_2(x)={}&y_{k-1}
\frac{(x-x_k)(x-x_{k+1})}
{(x_{k-1}-x_k)(x_{k-1}-x_{k+1})}\\
&+y_k
\frac{(x-x_{k-1})(x-x_{k+1})}
{(x_k-x_{k-1})(x_k-x_{k+1})}\\
&+y_{k+1}
\frac{(x-x_{k-1})(x-x_k)}
{(x_{k+1}-x_{k-1})(x_{k+1}-x_k)}.
\end{aligned}
\]

\subsubsection{2.2.3 Lagrange 插值多项式}\label{lagrange-ux63d2ux503cux591aux9879ux5f0f}

上面对 \(n=1\) 及 \(n=2\) 的情况，得到了一次与二次插值多项式 \(L_1(x)\) 及 \(L_2(x)\)，它们分别由式（2.2.5）和式（2.2.7）表示。这种用插值基函数表示的方法容易推广到一般情形。下面讨论通过 \(n+1\) 个节点 \(x_0<x_1<\cdots<x_n\) 的 \(n\) 次插值多项式 \(L_n(x)\)，假定它满足条件

\[
L_n(x_j)=y_j\qquad(j=0,1,\cdots,n).
\tag{2.2.8}
\]

为了构造 \(L_n(x)\)，先定义 \(n\) 次插值基函数。

\textbf{定义 2.2} 若 \(n\) 次多项式 \(l_j(x)\)（\(j=0,1,\cdots,n\)）在 \(n+1\) 个节点 \(x_0<x_1<\cdots<x_n\) 上满足条件

\[
l_j(x_k)=
\begin{cases}
1,&k=j,\\
0,&k\ne j
\end{cases}
\qquad(j,k=0,1,\cdots,n),
\tag{2.2.9}
\]

就称这 \(n+1\) 个 \(n\) 次多项式 \(l_0(x),l_1(x),\cdots,l_n(x)\) 为节点 \(x_0,x_1,\cdots,x_n\) 上的 \textbf{\(n\) 次插值基函数}。

\(n=1\) 及 \(n=2\) 时的情况前面已经讨论。用类似的推导方法，可得到 \(n\) 次插值基函数为

\[
l_k(x)=
\frac{(x-x_0)\cdots(x-x_{k-1})(x-x_{k+1})\cdots(x-x_n)}
{(x_k-x_0)\cdots(x_k-x_{k-1})(x_k-x_{k+1})\cdots(x_k-x_n)}
\quad(k=0,1,\cdots,n).
\tag{2.2.10}
\]

显然它满足条件（2.2.9）。于是，满足条件（2.2.8）的插值多项式 \(L_n(x)\) 可表示为

\[
L_n(x)=\sum_{k=0}^{n}y_kl_k(x).
\tag{2.2.11}
\]

由 \(l_k(x)\) 的定义知

\[
L_n(x_j)=\sum_{k=0}^{n}y_kl_k(x_j)=y_j
\qquad(j=0,1,\cdots,n).
\]

形如式（2.2.11）的插值多项式 \(L_n(x)\) 称为 \textbf{Lagrange 插值多项式}，而式（2.2.5）和式（2.2.7）是其 \(n=1\) 和 \(n=2\) 的特殊情形。若引入记号

\[
\omega_{n+1}(x)=(x-x_0)(x-x_1)\cdots(x-x_n),
\tag{2.2.12}
\]

容易求得

\[
\omega'_{n+1}(x_k)
=(x_k-x_0)\cdots(x_k-x_{k-1})(x_k-x_{k+1})\cdots(x_k-x_n),
\]

于是式（2.2.11）可改写成

\href{../../source-images/pdf-032.jpeg}{原页：书页 19／PDF 32}

\begin{quote}
页首接 PDF 31；页末线性插值余项续 PDF 33。
\end{quote}

\[
L_n(x)=\sum_{k=0}^{n}y_k
\frac{\omega_{n+1}(x)}
{(x-x_k)\omega'_{n+1}(x_k)}.
\tag{2.2.13}
\]

注意，\(n\) 次插值多项式 \(L_n(x)\) 通常是次数为 \(n\) 的多项式，特殊情况次数可能小于 \(n\)。例如，通过三点 \((x_0,\allowbreak y_0),\allowbreak (x_1,\allowbreak y_1),\allowbreak (x_2,\allowbreak y_2)\) 的二次插值多项式 \(L_2(x)\)，如果三点共线，则 \(y=L_2(x)\) 就是一直线，而不是抛物线，这时 \(L_2(x)\) 是一次式。

\subsubsection{2.2.4 插值余项}\label{ux63d2ux503cux4f59ux9879}

若在 \([a,b]\) 上用 \(L_n(x)\) 近似 \(f(x)\)，则其截断误差为 \(R_n(x)=f(x)-L_n(x)\)，\(R_n(x)\) 也称为插值多项式的余项或插值余项。关于插值余项估计有以下定理。

\textbf{定理 2.2} 设 \(f^{(n)}(x)\) 在 \([a,b]\) 上连续，\(f^{(n+1)}(x)\) 在 \((a,b)\) 内存在，节点 \(a\le x_0<x_1<\cdots<x_n\le b\)，\(L_n(x)\) 是满足条件式（2.2.8）的插值多项式，则对于任何 \(x\in[a,b]\)，插值余项

\[
R_n(x)=f(x)-L_n(x)=\frac{f^{(n+1)}(\xi)}{(n+1)!}\omega_{n+1}(x).
\tag{2.2.14}
\]

这里，\(\xi\in(a,b)\) 且依赖于 \(x\)，\(\omega_{n+1}(x)\) 是式（2.2.12）所定义的。

\textbf{证明} 由给定条件知，\(R_n(x)\) 在节点 \(x_k\)（\(k=0,1,\cdots,n\)）上为零，即

\[
R_n(x_k)=0\qquad(k=0,1,\cdots,n),
\]

于是

\[
R_n(x)=K(x)(x-x_0)(x-x_1)\cdots(x-x_n)=K(x)\omega_{n+1}(x),
\tag{2.2.15}
\]

其中 \(K(x)\) 是与 \(x\) 有关的待定函数。

现把 \(x\) 看成 \([a,b]\) 上一个固定点，作函数

\[
\varphi(t)=f(t)-L_n(t)-K(x)(t-x_0)(t-x_1)\cdots(t-x_n),
\]

根据插值条件及余项定义，可知 \(\varphi(t)\) 在点 \(x_0,x_1,\cdots,x_n\) 及 \(x\) 处均为零，故 \(\varphi(t)\) 在 \([a,b]\) 上有 \(n+2\) 个零点，根据 Rolle 定理，\(\varphi'(t)\) 在 \(\varphi(t)\) 的两个零点间至少有一个零点，故 \(\varphi'(t)\) 在 \([a,b]\) 上至少有 \(n+1\) 个零点。对 \(\varphi'(t)\) 再应用 Rolle 定理，可知 \(\varphi''(t)\) 在 \([a,b]\) 上至少有 \(n\) 个零点。依此类推，\(\varphi^{(n+1)}(t)\) 在 \((a,b)\) 内至少有一个零点，记作 \(\xi\in(a,b)\)，使

\[
\varphi^{(n+1)}(\xi)=f^{(n+1)}(\xi)-(n+1)!K(x)=0,
\]

于是

\[
K(x)=\frac{f^{(n+1)}(\xi)}{(n+1)!},\qquad \xi\in(a,b)
\]

且依赖于 \(x\)。

将上式代入式（2.2.15），就得到余项表达式（2.2.14）。证毕。

应当指出，余项表达式只有在 \(f(x)\) 的高阶导数存在时才能应用。\(\xi\) 在 \((a,b)\) 内的具体位置通常不可能给出，如果可以求出 \(\max_{a\le x\le b}|f^{(n+1)}(x)|=M_{n+1}\)，那么插值多项式 \(L_n(x)\) 逼近 \(f(x)\) 的截断误差限是

\[
|R_n(x)|\le\frac{M_{n+1}}{(n+1)!}|\omega_{n+1}(x)|.
\tag{2.2.16}
\]

当 \(n=1\) 时，线性插值余项为

\href{../../source-images/pdf-033.jpeg}{原页：书页 20／PDF 33}

\begin{quote}
页首续 PDF 32 的线性插值余项；末行误差估计的不等式续 PDF 34。
\end{quote}

\[
R_1(x)=\frac12 f''(\xi)\omega_2(x)
=\frac12 f''(\xi)(x-x_0)(x-x_1),
\qquad \xi\in[x_0,x_1];
\tag{2.2.17}
\]

当 \(n=2\) 时，抛物插值的余项为

\[
R_2(x)=\frac16 f'''(\xi)(x-x_0)(x-x_1)(x-x_2),
\qquad \xi\in[x_0,x_2].
\tag{2.2.18}
\]

\textbf{例 2.1} 已知 \(\sin0.32=0.314\,567\)，\(\sin0.34=0.333\,487\)，\(\sin0.36=0.352\,274\)，用线性插值及抛物插值计算 \(\sin0.336\,7\) 的值，并估计截断误差。

\textbf{解} 由题意取 \(x_0=0.32\)，\(y_0=0.314\,567\)，\(x_1=0.34\)，\(y_1=0.333\,487\)，\(x_2=0.36\)，\(y_2=0.352\,274\)。

用线性插值计算，取 \(x_0=0.32\) 及 \(x_1=0.34\)，由式（2.2.3）得

\[
\begin{aligned}
\sin0.336\,7\approx L_1(0.336\,7)
&=y_0+\frac{y_1-y_0}{x_1-x_0}(0.336\,7-x_0)\\
&=0.314\,567+\frac{0.018\,92}{0.02}\times0.016\,7\\
&=0.330\,365.
\end{aligned}
\]

其截断误差由式（2.2.17）得

\[
|R_1(x)|\le\frac{M_2}{2}|(x-x_0)(x-x_1)|,
\]

其中

\[
M_2=\max_{x_0\le x\le x_1}|f''(x)|.
\]

因

\[
f(x)=\sin x,\qquad f''(x)=-\sin x,
\]

可取

\[
M_2=\max_{x_0\le x\le x_1}|\sin x|=\sin x_1\le0.333\,5,
\]

于是

\[
\begin{aligned}
|R_1(0.336\,7)|
&=|\sin0.336\,7-L_1(0.336\,7)|\\
&\le\frac12\times0.333\,5\times0.016\,7\times0.003\,3\\
&\le0.92\times10^{-5}.
\end{aligned}
\]

若取 \(x_1=0.34\)，\(x_2=0.36\) 为节点，则线性插值为

\[
\begin{aligned}
\sin0.336\,7\approx\widetilde L_1(0.336\,7)
&=y_1+\frac{y_2-y_1}{x_2-x_1}(0.336\,7-x_1)\\
&=0.333\,487+\frac{0.018\,787}{0.02}\times(-0.003\,3)\\
&=0.330\,387,
\end{aligned}
\]

其截断误差为

\[
|\widetilde R_1(x)|\le\frac{M_2}{2}|(x-x_1)(x-x_2)|,
\]

其中

\[
M_2=\max_{x_0\le x\le x_2}|f''(x)|\le0.352\,3.
\]

于是

\[
|\widetilde R_1(0.336\,7)|=|\sin0.336\,7-\widetilde L_1(0.336\,7)|
\]

\href{../../source-images/pdf-034.jpeg}{原页：书页 21／PDF 34}

\begin{quote}
页首续 PDF 33 的 \(|\widetilde R_1(0.3367)|\) 误差估计。
\end{quote}

\[
\begin{aligned}
&\le\frac12\times0.352\,3\times0.003\,3\times0.023\,3\\
&\le1.36\times10^{-5}.
\end{aligned}
\]

用抛物插值计算 \(\sin0.336\,7\) 时，由式（2.2.7）得

\[
\begin{aligned}
\sin0.336\,7\approx{}&
y_0\frac{(x-x_1)(x-x_2)}{(x_0-x_1)(x_0-x_2)}
+y_1\frac{(x-x_0)(x-x_2)}{(x_1-x_0)(x_1-x_2)}\\
&+y_2\frac{(x-x_0)(x-x_1)}{(x_2-x_0)(x_2-x_1)}\\
={}&L_2(0.336\,7)\\
={}&0.314\,567\times\frac{0.768\,9\times10^{-4}}{0.000\,8}
+0.333\,487\times\frac{3.89\times10^{-4}}{0.000\,4}\\
&+0.352\,274\times\frac{-0.551\,1\times10^{-4}}{0.000\,8}\\
={}&0.330\,374.
\end{aligned}
\]

这个结果与六位有效数字的正弦函数表完全一样，这说明查表时用二次插值精度已相当高了。其截断误差限由式（2.2.18）得

\[
|R_2(x)|\le\frac{M_3}{6}|(x-x_0)(x-x_1)(x-x_2)|,
\]

其中

\[
M_3=\max_{x_0\le x\le x_2}|f'''(x)|=\cos x_0<0.828,
\]

于是

\[
\begin{aligned}
|R_2(0.336\,7)|
&=|\sin0.336\,7-L_2(0.336\,7)|\\
&\le\frac16\times0.828\times0.016\,7\times0.003\,3\times0.023\,3\\
&\le0.178\times10^{-7}.
\end{aligned}
\]

\begin{center}\rule{0.5\linewidth}{0.5pt}\end{center}

\subsection{相关原页校注（agent 补充）}\label{ux76f8ux5173ux539fux9875ux6821ux6ce8agent-ux8865ux5145}

以下校注来自本节覆盖的原页，可能也涉及同页相邻小节；原文照录与校正说明分开保留。

\subsubsection{原书 PDF 页 34 的校注}\label{ux539fux4e66-pdf-ux9875-34-ux7684ux6821ux6ce8}

\href{../pages/pdf-034.md}{完整原页转录} · \href{../../source-images/pdf-034.jpeg}{原页图像}

校注记录：NA5E-E002。

\subsubsection{校注（agent 补充）：NA5E-E002}\label{ux6821ux6ce8agent-ux8865ux5145na5e-e002}

上面的 \(0.828\)、\(0.178\times10^{-7}\) 和代入行的 \(3.89\times10^{-4}\) 均经过原页放大核对，确为扫描原书印字。\textbf{这几行不能直接作为正确的数值结果使用。}

\begin{enumerate}
\def\labelenumi{\arabic{enumi}.}
\tightlist
\item
  本例以弧度计算，\(\cos0.32\approx0.9492354180824409\)，所以 \(\cos x_0<0.828\) 不成立。可用保守上界 \(M_3<0.95\)。
\item
  即使照原书使用 \(0.828\)，该乘积也等于 \(1.77200694\times10^{-7}\)，大于原文所写的 \(0.178\times10^{-7}=1.78\times10^{-8}\)。
\item
  对正确的导数上界，用原节点得到的截断误差上界为
\end{enumerate}

\[
|R_2(0.3367)|
\le\frac{0.95}{6}\,(0.0167)(0.0033)(0.0233)
=2.03309975\times10^{-7}.
\]

这里的 \(R_2\) 指用精确节点函数值定义的插值截断误差。

\begin{enumerate}
\def\labelenumi{\arabic{enumi}.}
\setcounter{enumi}{3}
\tightlist
\item
  二次插值代入行的中间乘积应为 \(0.0167\times0.0233=0.00038911=3.8911\times10^{-4}\)。若严格按原书印出的 \(3.89\times10^{-4}\) 计算整行，会得到 \(0.3302826531125\)，不能得到末行的 \(0.330374\)。保留精确乘积，并使用题目给的有限精度表值，可得
\end{enumerate}

\[
\begin{aligned}
L_1(0.3367)&=0.3303652,\\
\widetilde L_1(0.3367)&=0.330387145,\\
L_2(0.3367)&=0.3303743620375.
\end{aligned}
\]

最后一项保留六位有效数字确实为 \(0.330374\)，与 \(\sin0.3367\approx0.3303741915556280\) 的六位有效数字一致。

\subsection{本节覆盖原页的校注记录}\label{ux672cux8282ux8986ux76d6ux539fux9875ux7684ux6821ux6ce8ux8bb0ux5f55}

下列记录按原页关联，含同页相邻小节的校注；计算时同时读取上文校注和完整原页。

\begin{itemize}
\tightlist
\item
  \texttt{NA5E-E002}：原书 PDF 页 34
\end{itemize}

\end{document}
